\section{问题重述}
\subsection{问题背景}
研究对象是一个通用神经网络处理器，在多个核心上执行同一张计算图。矩阵、向量及搬运操作之间有先后依赖；每核片上存储有限，所有核心共享外部存储带宽。切图既释放并行性，也改变跨核搬运、等待和张量驻留时间，因此需要同时确定操作归属与执行顺序，并以总体任务执行时间Makespan（下文亦称工期）评价方案。

此类问题涉及多个研究方向。在空间加速器的数据流映射上，Eyeriss把矩阵运算的行固定数据流写成静态排列，每条计算链就地复用片上缓冲，以减少外部存储访问\cite{chen2016eyeriss}；Chen进一步系统地分析了如何把数据流选择与存储层次设计联合考虑，归纳出灵活适应不同网络的加速器架构原则\cite{chen2018phd}。在多芯片与资源划分上，Simba把多核推理拆分为多个片上小核模块，按模块间路由代价分配算子，以减少片间搬运瓶颈\cite{shao2019simba}；FlexFlow探索了不同数据流配置对片上带宽的敏感性，总结出可在硬件参数变化时保持效率的灵活映射接口\cite{lu2017flexflow}。在计算图的放置与调度上，Mirhoseini等用强化学习优化深度学习模型在异构设备上的算子放置，以端到端工期为优化目标\cite{mirhoseini2017placement}。在数据复用代价建模上，Kwon等建立了统一的数据流代价框架，分析不同映射下的存储访问量与带宽压力之间的权衡关系\cite{kwon2019dataflow}。在生产者—消费者管道的存储调度上，Ragan-Kelley通过分离算法与调度决策，推导出Halide式流水线中的最优生产次序与缓冲大小\cite{ragankelley2014phd}。上述工作覆盖了从单核数据流到多核图划分、从静态映射到搜索型调度的多个层面；本题硬件结构和执行规则较固定，三问沿用同一基准，逐步扩展可复用范围。

三问的实验对象是同一批100张计算图，配合固定的硬件配置参数，由同一个事件评估程序评价。提交内容为全部非COPY操作的子图归属及各核的子图顺序；评估程序据此生成搬运与缓存换出换入操作，计算Makespan、额外数据搬运量和缓存命中率。L1和UB由各核独占，DDR由多核共享，问题三另设共享只读L2 Cache。

\subsection{问题描述}
在相同计算图和固定硬件配置下，三问依次改变Task组织、片上复用和共享读服务条件。各问均须给出可执行的多核计划及相应评价指标。

\textbf{问题一：}场景A将每个子图构造成独立Task，跨子图张量通过DDR中转，同核相邻Task也不保留片上副本。需要确定全部非COPY操作的子图归属和每核的子图执行次序，在满足计算依赖与Task等待规则的前提下缩短总体工期、控制额外搬运，并给出1至5核的加速比及逐图结果。

\textbf{问题二：}场景B把同核子图合并为核级Task，使同核依赖可以复用L1和UB中的数据；跨核依赖仍需写出、读入和固定同步延迟。需要在容量约束下重新设计切图、分核和核内顺序，给出不同核数的工期、加速比与额外搬运，并分析同核复用带来的收益和换出恢复的代价。

\textbf{问题三：}在问题二的核级Task结构上增加共享只读L2 Cache，本文将此配置记为C。读请求在发射时查询，完成后尝试插入，空间不足时按FIFO淘汰。需比较同图同核数下有无L2的Makespan与额外数据搬运量，给出命中率，并设计适应共享读资源的方案。
